\section{Related Work}
Typestate verification was first introduced by Strom and Yemini to track the initialization state of variables. Since then, many researchers have extended typestate analysis to object-oriented programming. Plural, developed by Bierhoff and Aldrich, uses fractional permissions to verify protocol compliance in Java-like languages. While Plural supports complex aliasing patterns, its annotation burden is substantial. In contrast, Synapse prioritizes simplicity by using a strict linear type system for actor references, avoiding the complexity of fractional permissions.

Session types, pioneered by Honda et al., provide another formalism for verifying communication protocols. Session types specify the sequence and types of messages exchanged between two endpoints. While session types are expressive, integrating them into general-purpose object-oriented systems often requires complex compiler machinery. Synapse bridges this gap by encoding actor behaviors as state transitions in a conventional typestate framework. This approach allows developers to specify protocols using state transition tables directly inside class definitions, providing a more intuitive programming model. Recent work on the Vertex virtual machine by Nimbus and Pendelton~\cite{nimbus2019vertex} has explored high-performance execution of linear code; Synapse builds on this platform to offer both safety and execution efficiency.
